% Experiment Template
% 1) Copy this file and rename it, e.g. experiments/experiment_04_firewall_rules.tex
% 2) In main.tex, add:
%    \includeexperiment{XX}{Experiment Title}{experiments/your_new_file_name}
% 3) Keep this file content as sections only (no \experiment command here)

\section{Theory}
The Rail Fence cipher is a simple keyless transposition cipher that writes the plaintext in a zigzag pattern across a fixed number of rows, then reads off each row to produce the ciphertext. For example, encrypting \textbf{"HELLO WORLD"} with 2 rails yields \textbf{"HLOOLELWRD"}.

\section{Algorithm}
The encryption algorithm for rail fence cipher is,
\begin{enumerate}
  \item set ans = empty string
  \item n = length of cipher text
  \item for i = 0 to n-1 do
  \begin{enumerate}
    \item if i is even then ans = ans + cipher[i/2]
  \end{enumerate}
  \item for i = 0 to n-1 do
  \begin{enumerate}
    \item if i is odd then ans = ans + cipher[n/2 + i/2]
  \end{enumerate}
  \item return ans
\end{enumerate}
The decryption algorithm for rail fence cipher is,
\begin{enumerate}
  \item set ans = cipher text
  \item n = length of cipher text
  \item for i = 0 to n-1 do
  \begin{enumerate}
    \item if i is even then ans[i] = cipher[i/2]
    \item if i is odd then ans[i] = cipher[n/2 + i/2]
  \end{enumerate}
  \item return ans
\end{enumerate}


% \section{Flowchart (motherboard)}
% \begin{figure}[H]
%   \centering
%   \begin{autofitflowchart}
%     \node (start) [terminator] {Start};
%     \node (input) [io, below=of start] {Input};
%     \node (process1) [process, below=of input] {Process};
%     \node (decision1) [decision, below=of process1, yshift=-0.2cm] {Condition?};
%     \node (loopstep) [process, right=of decision1, xshift=1.2cm] {Update / Loop};
%     \node (output) [io, below=of decision1, yshift=-0.3cm] {Output};
%     \node (end) [terminator, below=of output] {End};

%     \draw [line] (start) -- (input);
%     \draw [line] (input) -- (process1);
%     \draw [line] (process1) -- (decision1);

%     \draw [line] (decision1.east) -- node[above]{No} (loopstep.west);
%     \draw [line] (loopstep.north) |- ([xshift=1.0cm]process1.east) |- (process1.east);

%     \draw [line] (decision1) -- node[right]{Yes} (output);
%     \draw [line] (output) -- (end);
%   \end{autofitflowchart}
%   \caption{Flowchart title}
% \end{figure}

\section{Code}
\begin{minted}{cpp}
#include<bits/stdc++.h>
using namespace std;

class rail {
  public:
    rail();
    string encrypt(string s);
    string decrypt(string s);
};

rail::rail() {
}

string rail::encrypt(string s) {
  string ans;
  int n = s.length();
  if (n%2 == 1) {
    s+=' ';
    n++;
  }
  for (int i = 0; i < n; i+=2) {
    ans+= s[i];
  }
  for (int i = 1; i < n; i+=2) {
    ans += s[i];
  }
  return ans;
}

string rail::decrypt(string s) {
  string ans = s;
  int n = s.length();
  for (int i = 0; i < n; i++) {
    if (i%2 == 0) {
      ans[i] = s[i/2];
    } else {
      ans[i] = s[n/2 + i/2];
    }
  }
  return ans;
}

int main() {
  rail cipher;
  cout << "Enter Plaintext : ";
  string text;
  getline(cin, text);
  string encrypted = cipher.encrypt(text);
  cout << "Encrypted text: " << encrypted << endl;
  string decrypted = cipher.decrypt(encrypted);
  cout << "Decrypted text: " << decrypted << endl;
  return 0;
}

\end{minted}

\section{Result and Discussion - I/O}
\begin{minted}{text}
Enter Plaintext : This is password
Encrypted text: Ti spswrhsi asod
Decrypted text: This is password
\end{minted}


\section{Conclusion}
Rail fence cipher is a simple keyless transposition cipher that is easy to encrypt, decrypt, and break, making it a useful introduction to transposition ciphers.
